\documentclass[a4paper,11pt]{article}
\usepackage[T2A]{fontenc}
\usepackage[utf8]{inputenc}
\usepackage[russian]{babel}
\usepackage{amsmath,amssymb}
\usepackage{geometry}
\geometry{left=21mm,right=21mm,top=19mm,bottom=19mm,headheight=23pt,headsep=6mm,footskip=13mm}
\usepackage{microtype}
\usepackage{xcolor}
\usepackage{parskip}
\usepackage{fancyhdr}
\usepackage{tikz}
\usepackage{setspace}
\usepackage{hyperref}
\definecolor{accent}{HTML}{006F72}
\definecolor{darkaccent}{HTML}{004D4D}
\definecolor{muted}{HTML}{545F62}
\definecolor{rulecolor}{HTML}{9CBEBF}
\hypersetup{colorlinks=true,linkcolor=darkaccent,urlcolor=darkaccent}
\setlength{\parindent}{0pt}
\setlength{\parskip}{2pt}
\setstretch{1.00}
\setlength{\abovedisplayskip}{5pt}
\setlength{\belowdisplayskip}{5pt}
\setlength{\abovedisplayshortskip}{3pt}
\setlength{\belowdisplayshortskip}{3pt}
\pagestyle{fancy}
\fancyhf{}
\fancyhead[L]{\small\color{muted}Проверка домашней работы}
\fancyhead[R]{\small\color{muted}Ярослав Верещагин}
\renewcommand{\headrulewidth}{0.5pt}
\renewcommand{\headrule}{\hbox to\headwidth{\color{rulecolor}\leaders\hrule height \headrulewidth\hfill}}
\fancyfoot[C]{\footnotesize\color{muted}Лёвин Артём Александрович эксклюзивно для Ярослава Верещагина}
\renewcommand{\footrulewidth}{0pt}
\newcommand{\accentsection}[1]{\vspace{1pt}{\large\bfseries\color{darkaccent}#1}\par\vspace{0pt}{\color{rulecolor}\rule{\linewidth}{0.5pt}}\vspace{1pt}}
\begin{document}
\thispagestyle{empty}
\begin{center}
\vspace*{23mm}
{\large\color{muted}МАТЕМАТИКА}\par
\vspace{18mm}
{\Huge\bfseries\color{darkaccent}Проверка\par домашней работы}\par
\vspace{9mm}
{\Large Алгебраические преобразования\par и текстовые задачи}\par
\vspace{17mm}
{\color{accent}\rule{0.68\linewidth}{0.8pt}}\par
\vspace{13mm}
{\large Ярослав Верещагин}\par
\vspace{4mm}
{\normalsize 23 ноября 2026 г.}\par
\end{center}
\vfill
\begin{center}
{\small\color{muted}Лёвин Артём Александрович\\
эксклюзивно для Ярослава Верещагина}
\end{center}
\vspace{11mm}
\begin{tikzpicture}[remember picture,overlay]
 \draw[accent!55,line width=0.8pt] ([xshift=22mm,yshift=24mm]current page.south west)
 .. controls ([xshift=74mm,yshift=42mm]current page.south west)
 and ([xshift=-73mm,yshift=11mm]current page.south east)
 .. ([xshift=-22mm,yshift=25mm]current page.south east);
\end{tikzpicture}
\clearpage

\section*{\color{accent}Задание №\,8. Наборщик}
\noindent\textbf{Условие.} Наборщик должен набрать 240 страниц. Ежедневно печатая на 4 страницы больше плана, он закончил работу на 2 дня раньше. Сколько страниц в день было запланировано?

\accentsection{Где возникла ошибка}
Уравнение составлено верно. Ошибка появилась при делении его на два после избавления от знаменателей.

До ошибочного шага получено правильное равенство:
\[
240x+240\cdot4-240x-2(x^2+4x)=0.
\]
В следующей строке произведение \(240\cdot4\) было заменено после деления на два на \(120\cdot2\). Но это означает, что \emph{оба} множителя разделили на два, то есть всё произведение уменьшили в четыре раза вместо двух:
\[
\underbrace{\frac{240\cdot4}{2}}_{\text{верно}}=120\cdot4=480,
\qquad
\underbrace{120\cdot2}_{\text{ошибка}}=240.
\]
Из-за этого получено неверное квадратное уравнение и неверный дискриминант. При делении произведения на число достаточно разделить \textbf{один} из множителей.

\accentsection{Полное правильное решение}
Пусть \(x\)~--- запланированное число страниц в день, \(x>0\). По факту наборщик печатал \(x+4\) страницы в день. На работу по плану требуется \(240/x\) дней, а по факту~--- \(240/(x+4)\) дней. Разница составляет два дня:
\[
\frac{240}{x}-\frac{240}{x+4}=2.
\]
Приведём дроби к общему знаменателю и раскроем скобки:
\begin{align*}
\frac{240(x+4)-240x}{x(x+4)}&=2,\\
\frac{240x+960-240x}{x(x+4)}&=2,\\
\frac{960}{x(x+4)}&=2.
\end{align*}
Так как \(x>0\), знаменатель не равен нулю. Умножим обе части на \(x(x+4)\):
\[
960=2x(x+4)\quad\Longleftrightarrow\quad
480=x^2+4x\quad\Longleftrightarrow\quad x^2+4x-480=0.
\]
Найдём дискриминант и корни:
\[
D=4^2-4\cdot1\cdot(-480)=1936=44^2,
\qquad x=\frac{-4\pm44}{2}.
\]
Получаем \(x_1=20\) и \(x_2=-24\). Отрицательное количество страниц в день невозможно, поэтому \(x=20\).

\textbf{Проверка.} По плану нужно \(240:20=12\) дней; при наборе по \(24\) страницы~--- \(240:24=10\) дней. Работа действительно закончена на два дня раньше.

\textbf{Ответ:} \(20\) страниц в день.

\accentsection{Задача на закрепление}
Наборщику нужно набрать 360 страниц. Если ежедневно набирать на 6 страниц больше плана, работа завершится на два дня раньше. Сколько страниц в день он планировал набирать?

\end{document}
